% Part: computability
% Chapter: computability-theory
% Section: applications-fixed-point

\documentclass[../../../include/open-logic-section]{subfiles}

\begin{document}

\olfileid{cmp}{thy}{apf}
\olsection{د ثابت ټکي د قضيې کارول}

د ثابت ټکي قضيه په اصل کښې اجازه راکوي چې جزوي محاسبه کېدونکې تابعې
د هغو د شاخصونو په وسيله تعريف کړو۔ د بېلګې په توګه، داسې شاخص~$e$
موندلے شو چې د هر~$y$ دپاره
\[
\cfind{e}(y) = e + y.
\]
د بلې بېلګې په توګه، د ثابت ټکي د قضيې ثبوت کارولے شو څو په جاوا يا
C++ کښې داسې پروګرام جوړ کړو چې خپل متن چاپ کړي۔

راياد کړئ چې که د هر~$e$ دپاره~$W_e$ د~$\cfind{e}$ د تعريف ساحه
وټاکو، نو لړۍ~$W_0$، $W_1$، $W_2$،~\dots په محاسبوي ډول د شمېر وړ
سټونه شمېري۔ ځينې دا سټونه محاسبه کېدونکي دي۔ پوښتلے شو چې ايا داسې
الګوريتم شته چې قيمت~$x$ د ننوت په توګه واخلي، او که~$W_x$ محاسبه
کېدونکے وي، د هغه د ځانګړونکې تابع يو شاخص را وګرځوي؟ ځواب ``نه''
دے، داسې الګوريتم نشته:

\begin{thm}
داسې جزوي محاسبه کېدونکې تابع~$f$ نشته چې دا خاصيت ولري: هر کله چې
~$W_e$ محاسبه کېدونکے وي، نو~$f(e)$ تعريف شوے وي او~$\cfind{f(e)}$
ئې ځانګړونکې تابع وي۔
\end{thm}

\begin{proof}
$f$ هره جزوي محاسبه کېدونکې تابع وي؛ داسې~$e$ به جوړ کړو چې~$W_e$
محاسبه کېدونکے وي، خو~$\cfind{f(e)}$ ئې ځانګړونکې تابع نۀ وي۔ د ثابت
ټکي د قضيې په کارولو داسې شاخص~$e$ موندلے شو چې
\[
\cfind{e}(y) \simeq
\begin{cases}
0 & \text{که $y=0$ او $\cfind{f(e)}(0) \fdefined = 0$ وي} \\
\text{تعريف شوې نۀ ده} & \text{که داسې نۀ وي۔}
\end{cases}
\]
يعنې~$e$ د ثابت ټکي د قضيې په دې تابع له پلي کولو تر لاسه کېږي:
\[
g(x,y) \simeq
\begin{cases}
0 & \text{که $y=0$ او $\cfind{f(x)}(0) \fdefined = 0$ وي} \\
\text{تعريف شوې نۀ ده} & \text{که داسې نۀ وي۔}
\end{cases}
\]
په غير رسمي ډول داسې ليدلے شو چې~$g$ جزوي محاسبه کېدونکې ده: پر
ننوتونو~$x$ او~$y$ الګوريتم لومړے ګوري چې~$y$ له~$0$ سره برابر دے
که نه۔ که برابر وي، الګوريتم~$f(x)$ محاسبه کوي، او بيا نړيوال ماشين
کاروي څو~$\cfind{f(x)}(0)$ محاسبه کړي۔ که دا وروستۍ محاسبه ودرېږي او
~$0$ را وګرځوي، الګوريتم~$0$ راګرځوي؛ که نه، الګوريتم نۀ درېږي۔
\textbf{د ژباړې د نسخې سمون (OLCMP-036):} سرچينې د ثبوت په پيل کښې
خپل اختياري سيال بشپړه محاسبه کېدونکې تابع بللې وه، حال دا چې د
قضيې دعوه د ټولو جزوي محاسبه کېدونکو تابعو په اړه ده۔ دلته د قضيې
هماغه جزوي ټولګے ساتل شوے؛ ورپسې جوړښت د ناتعريف حالت په ګډون هم
کار کوي۔

خو اوس پام وکړئ چې که~$\cfind{f(e)}(0)$ تعريف شوې او له~$0$ سره
برابره وي، نو~$\cfind{e}(y)$ کټ مټ هغه وخت تعريف شوې ده چې~$y$ له
~$0$ سره برابر وي، نو $W_e = \{ 0 \}$۔ که~$\cfind{f(e)}(0)$ تعريف
شوې نۀ وي، يا تعريف شوې خو له~$0$ سره برابره نۀ وي، نو
$W_e = \emptyset$۔ په دواړو حالاتو کښې~$\cfind{f(e)}$ د~$W_e$
ځانګړونکې تابع نۀ ده، ځکه پر ننوت~$0$ ناسم ځواب ورکوي۔
\end{proof}

\end{document}
